\scanpage{14}\mspage{11}

\[
 \left\{\begin{aligned}
 L\Lambda&=1-\sum_{0\le i\le n}\pi_{i,0}\\
 \Lambda L&=1-\sum_{n\le j\le 2n}\pi_{\deleted{\unsure{2i}}j,0}
 \end{aligned}\right.
 \tag*{(6.15)}
\]
\deleted{Quant aux $\pi_{j,0}$, ils sont déterminés par les}

\uline{Question} Donner un système de \deleted{\unsure{générateurs}} relations en les $\Lambda,L$ exclusivement, \deleted{\unsure{puis}} des relations exprimant les $\pi_i$ comme polynômes non comm en $L,\Lambda$ ou qui définissent $\Lambda$ en termes de $L$ et des $\pi_i$ \ldots
% slanted
\margin{cf n\textsuperscript{o} 7 pour l'anneau $\Phi$ : système de relations \ldots}

7. Revenant à un anneau gradué abstrait $\unsure{E|\mathcal E}$, muni d'un
\[
 L\in E^{(2)},
\]
on veut abstraire la situation du n\textsuperscript{o} précédent, sans utiliser $M^*$, en termes des propriétés du seul anneau \deleted{\unsure{$E$}} gradué $E^*$ muni de $L\in E^{(2)}$.

Partons pour ceci d\deleted{\unsure{e}}'un module libre $\mathbb{M}_{\unsure{0|\circ}}$ de rang $2n$ sur $\mathbb Z$, muni d'une 2-forme alternée non dégénérée \inserted{(sur $\mathbb Z$ !)} $\unsure{\xi|\zeta}_0\in\bigwedge^2\mathbb M_0$, \inserted{$\xi^*\in\bigwedge^2\mathbb M_0^\vee$ la forme duale,} soit \deleted{\unsure{$\xi$}} $\mathbb M_0^*=\bigwedge\mathbb M_0$, et considérons sur $M_0^\alpha$ les opérateurs
\[
 \mathbb L x=\xi_0x,\qquad \mathbb{\Lambda}_0x=\xi_0^*\mathbin{\lrcorner}x .
\]
et considérons l'algèbre
\[
 \Phi_0=\mathbb Z[\mathbb L_0,\mathbb{\Lambda}_0]\subset\operatorname{End}_{\mathbb Z}(M_0^*)\qquad\text{anneau \unsure{engendré|gradué}.}
\]
Noter que $\operatorname{End}(\mathbb M_0^*)$, \inserted{opérant sur $\mathbb M_0^*$,} muni de $\mathbb L$, satisfait à la condition de (6.1) \unsure{(N)|($\forall$)|($\omega$)}, \deleted{\unsure{donc}} et \deleted{\unsure{$\xi$}} $\mathbb{\Lambda}_0$ est l'opération associée définie par (6.5).
% archive sheet number: 13
